\chapter{Materials and Methods}
\label{ch:methods}

This chapter describes the design and implementation of the proposed receptionist system.
\secref{sec:hardware} introduces the hardware platform, including the Pepper robot and the compute infrastructure.
\secref{sec:architecture} presents the overall system architecture and the communication model that binds its components.
Sections~\ref{sec:voice_pipeline} through~\ref{sec:orchestration} detail the individual subsystems: the voice interaction pipeline, knowledge grounding, non-verbal behavior, and session orchestration.
Finally, \secref{sec:implementation} summarizes the software stack and deployment strategy.
The complete source code is publicly available in the project repository\footnote{Repository URL to be added upon publication.}.

% ============================================================
\section{Hardware Platform}
\label{sec:hardware}

\subsection{Pepper Humanoid Robot}
\label{sec:pepper_robot}

Pepper is a semi-humanoid robot manufactured by SoftBank Robotics, designed for social interaction in public environments.
The robot stands approximately 1.2\,m tall and features 20 degrees of freedom, including a mobile base with three omnidirectional wheels, two articulated arms with five-fingered hands, and a head capable of pitch and yaw rotation.

For the purposes of this work, the following hardware capabilities are relevant:

\begin{itemize}
    \item \textbf{Microphones:} Four directional microphones mounted on the head, used natively for sound localization.
    In the proposed system, however, these are not used for speech recognition; audio input is captured externally (see \secref{sec:audio_delivery}).
    \item \textbf{Speakers:} Two speakers in the ears, accessible through the NAOqi audio service for programmatic playback.
    \item \textbf{Tablet:} A 10.1-inch touchscreen mounted on the chest, capable of displaying HTML content via a dedicated tablet service.
    \item \textbf{LEDs:} Multicolor LEDs around the eyes and on the shoulders, programmable for visual feedback.
    \item \textbf{Actuators:} Joint motors supporting predefined animation sequences, executed through the NAOqi animation player service.
\end{itemize}

Pepper's onboard computer (Intel Atom E3845, 4\,GB RAM) runs NAOqi OS and exposes a service-oriented API through the \texttt{qi} framework.
All programmatic access is performed over a TCP connection using the \texttt{libqi} Python bindings.
Critically, the onboard hardware is insufficient for running large language models or real-time speech recognition; all inference is offloaded to external compute, and the robot serves exclusively as a sensor and actuator platform.

\subsection{Compute Infrastructure}
\label{sec:compute_infra}

The system runs on a two-node compute architecture connected to the robot over a local network:

\begin{enumerate}
    \item \textbf{Raspberry Pi 5 (8\,GB):} Serves as the primary edge controller, co-located with Pepper on the same LAN segment.
    It runs all containerized services --- LiveKit server, session manager, robot bridge, voice agent in cloud mode, and supporting infrastructure --- and communicates with Pepper over TCP.

    \item \textbf{GPU server:} A departmental server equipped with NVIDIA GPUs, accessible via SSH through a university jump host.
    It hosts the local-mode voice agent and the vLLM inference server (see \secref{sec:local_mode}).
    Because the server is not directly reachable from the RPi, an autossh-based reverse SSH tunnel forwards the necessary services so that the GPU server can reach RPi-hosted components transparently.
\end{enumerate}

% TODO: \figref{fig:network_topology} — network topology diagram showing RPi, Pepper, jump host, GPU server, and the tunneled connections.

The choice to split compute between an edge device and a remote GPU server reflects a key design constraint: the cloud-mode agent benefits from low-latency co-location with the LiveKit server on the RPi, while the local-mode agent requires GPU resources that only the departmental server can provide.
This two-node arrangement enables a fair comparison of both deployment paradigms without hardware changes.


% ============================================================
\section{System Architecture}
\label{sec:architecture}

\subsection{High-Level Architecture}
\label{sec:high_level_arch}

The system is decomposed into loosely coupled services, each responsible for a distinct function.
This microservice-oriented design was motivated by three practical requirements:

\begin{enumerate}
    \item \textbf{Runtime isolation:} The robot bridge requires the \texttt{qi} Python bindings, which are compiled for a specific Python version and depend on platform-specific native libraries.
    Isolating this component in its own container avoids dependency conflicts with the voice agent and other services.

    \item \textbf{Independent restartability:} During development and live demonstrations, individual components (e.g., the voice agent after a prompt change, or the bridge after a Pepper reboot) must be restartable without disrupting the rest of the system.
    A monolithic process would require full system restarts for any change.

    \item \textbf{Distributed deployment:} The local-mode voice agent runs on a different physical machine (GPU server) than the remaining services (RPi).
    A service-based architecture naturally accommodates this split without code changes.
\end{enumerate}

The system consists of the following principal services:

\begin{itemize}
    \item \textbf{Voice Agent} — performs dialogue reasoning, decides tool calls, and generates spoken responses.
    Two variants exist: a cloud-mode agent running on the RPi and a local-mode agent running on the GPU server.
    Both register with LiveKit under unique names so the session manager can dispatch to the correct one.

    \item \textbf{Session Manager} — the central orchestrator.
    Creates LiveKit rooms, manages session state, dispatches voice agents by mode, tracks component health, and serves an Operator Panel web interface for monitoring and control.

    \item \textbf{Robot Bridge} — translates high-level commands (animation triggers, tablet content, audio streams) into NAOqi service calls on Pepper.

    \item \textbf{Audio Bridge} — joins the LiveKit room as a participant, subscribes to the voice agent's audio track, resamples the PCM stream, and forwards it to the robot bridge for playback on Pepper's speakers.

    \item \textbf{Room Monitor} — observes LiveKit room state and forwards transcripts to the session manager for display in the Operator Panel.

    \item \textbf{Infrastructure services} — LiveKit Server (WebRTC signaling and media routing), Redis (LiveKit state backend), and Weaviate (vector database for RAG).
\end{itemize}

% TODO: \figref{fig:architecture} — architecture diagram showing all services, their host locations (RPi vs. GPU server), and connections.

\subsection{Communication Model}
\label{sec:communication_model}

Three distinct communication channels bind the services together:

\begin{enumerate}
    \item \textbf{WebRTC via LiveKit.}
    All real-time audio transport passes through LiveKit rooms.
    The user (browser client or RPi microphone) publishes audio into a room; the voice agent subscribes to it, processes it, and publishes its response audio back.
    The audio bridge subscribes to the agent's output track and extracts the PCM stream for Pepper.
    LiveKit was chosen for its built-in acoustic echo cancellation, codec handling, room-level participant management, and support for data channels used for session-control signaling.

    \item \textbf{HTTP REST.}
    Tool calls from the voice agent to external services use synchronous HTTP requests.
    Animation triggers are sent to the robot bridge, which validates the request, dispatches the animation to Pepper in a background thread, and returns an immediate acknowledgement.
    Knowledge queries are routed to Weaviate's search API.
    Various services also report events (transcripts, metrics, tool calls) to the session manager via HTTP.

    \item \textbf{Raw TCP.}
    Audio playback on Pepper uses a custom length-prefixed TCP protocol between the audio bridge and the robot bridge.
    The audio bridge sends 16-bit PCM frames with a 4-byte length header; the robot bridge buffers and plays them through NAOqi's audio device service.
    Special sentinel values signal flush (end of utterance) and keepalive.
\end{enumerate}

This separation allows each channel to be optimized independently.
WebRTC handles the latency-sensitive bidirectional audio path with built-in reliability and codec negotiation.
HTTP provides a simple request-response interface for tool calls where sub-second latency is acceptable.
Raw TCP minimizes overhead for the high-throughput, unidirectional audio stream to the robot.


% ============================================================
\section{Voice Interaction Pipeline}
\label{sec:voice_pipeline}

The system supports two voice interaction modes that can be switched at runtime via the Operator Panel or a REST API call.
Both modes share the same tool set, session lifecycle, and robot control interface; they differ only in how speech is processed and how language reasoning is performed.

\subsection{Cloud Mode: OpenAI Realtime API}
\label{sec:cloud_mode}

The cloud-mode agent uses the OpenAI Realtime API with the \texttt{gpt-4o-mini-realtime} model.
This is an instance of the integrated end-to-end audio architecture described in \secref{sec:multimodal_voice}: the API accepts audio input directly, performs speech recognition, language reasoning, and speech synthesis within a single multimodal model, and streams synthesized audio back to the caller.

The agent is implemented as a LiveKit Agents SDK worker.
On startup, it creates an agent session with a single Realtime model instance that handles the entire speech-to-speech pipeline.
The model receives system instructions (persona definition, tool usage rules, animation guidelines) and two tool definitions (\texttt{query\_search} and \texttt{play\_animation}) that it can invoke during dialogue.

Because the Realtime API handles all speech processing internally, the agent process is lightweight and runs on the RPi without GPU requirements.
Co-locating it with the LiveKit server eliminates network latency between the agent and the WebRTC media path, which is critical for responsive turn-taking.
The only external network dependency is the OpenAI API itself, which introduces variable latency depending on server load.

\subsection{Local Mode: Cascaded STT--LLM--TTS}
\label{sec:local_mode}

The local-mode agent implements the modular cascaded speech architecture from \secref{sec:multimodal_voice}, consisting of three independent stages:

\begin{enumerate}
    \item \textbf{Speech-to-Text:} Faster Whisper, an optimized CTranslate2 port of OpenAI's Whisper model.
    The \texttt{tiny} model variant is used by default with INT8 quantization on CPU to minimize latency.
    The implementation includes silence detection based on RMS energy thresholds and minimum word count filtering to suppress hallucinated transcripts on quiet audio.

    \item \textbf{Language Model:} Qwen 2.5 7B Instruct, served via vLLM with an OpenAI-compatible API endpoint.
    The model runs with low temperature, sequential tool calling (no parallel calls), and thinking mode disabled.
    Tool calling uses the Hermes format, which vLLM's parser extracts and converts to structured function calls.

    \item \textbf{Text-to-Speech:} Piper, a fast neural TTS engine based on VITS.
    An English female voice ONNX model is used, providing natural-sounding speech synthesis on CPU with configurable speaking rate and prosody parameters.
\end{enumerate}

Voice activity detection is handled by Silero VAD, which determines when the user has started and stopped speaking, enabling the agent to process complete utterances rather than audio fragments.

The local-mode agent runs on the GPU server, where vLLM utilizes GPU resources for LLM inference.
It connects to the LiveKit server on the RPi through the SSH reverse tunnel, using LiveKit's built-in TURN relay (TLS over TCP) to transport WebRTC media, since SSH does not support UDP forwarding natively.

\paragraph{Practical Challenges with Local Tool Calling.}
Deploying structured tool calling with a 7B-parameter model revealed several engineering challenges that required application-level workarounds:

\begin{itemize}
    \item \textbf{Malformed JSON arguments:} Qwen 2.5 7B occasionally generates extra trailing braces or appends natural language text after the JSON object in tool call arguments.
    A sanitization layer intercepts function call arguments before execution, extracting the first balanced JSON object and discarding trailing content.
    The same sanitization is applied to the conversation history before each LLM call, preventing the inference server from crashing when it re-parses malformed arguments.

    \item \textbf{Tool description framing:} The model reliably calls tools it perceives as returning \emph{needed data} but tends to skip tools it views as fire-and-forget side effects.
    The animation tool description was rewritten to frame it as returning robot body state information that the model needs before speaking, rather than as a one-way gesture trigger.
    This increased call reliability from approximately 60\% to near-consistent invocation.

    \item \textbf{System prompt format sensitivity:} Multi-line system prompts with structured formatting caused the tool call parser to fail when the model interleaved natural language text with tool call markup.
    Condensing the system prompt to a single line resolved these parsing failures.
\end{itemize}

These workarounds are specific to the 7B model scale and may not be necessary with larger models.
They are documented here because they represent practical integration challenges that are not captured by standard benchmark evaluations of tool-calling capabilities.


% ============================================================
\section{Knowledge Grounding}
\label{sec:knowledge}

To answer factual questions about the faculty, the system implements a retrieval-augmented generation pipeline as described theoretically in \secref{sec:rag}.
This section details the concrete implementation.

\subsection{Document Ingestion and Embedding}
\label{sec:doc_ingestion}

The knowledge base consists of institutional documents relevant to the receptionist scenario: the CTU Statute, the FEE Statute, the Scholarship and Bursary Code, the Accommodation Code, the Career Code, and the Disciplinary Code.
These documents are stored as plain text files in the project repository.

On agent startup, a seeding routine checks whether the vector database collection exists.
If the collection is newly created, each document is inserted as an object containing its title, full text content, source reference, and an insertion timestamp.
Embedding is performed automatically by Weaviate using the \texttt{text-embedding-3-large} model from OpenAI.
Both the title and content fields are vectorized, and the resulting embeddings are stored alongside the original text for hybrid retrieval.

\subsection{Retrieval-Augmented Generation via Tool Calling}
\label{sec:rag_tool}

Knowledge retrieval is exposed to the LLM as a tool called \texttt{query\_search}.
When the model determines that answering a user question requires factual information, it emits a structured tool call with the user's query as the argument.
The tool execution follows these steps:

\begin{enumerate}
    \item The query string is passed to Weaviate's hybrid search, which combines BM25 keyword matching with vector similarity using a weighting parameter $\alpha = 0.7$, favoring semantic similarity over exact keyword matches.
    \item The top results are returned, each containing the document title, content, source, and a relevance score.
    \item The results are serialized as JSON and appended to the conversation context as a tool response.
    \item The LLM generates a natural-language answer grounded in the retrieved content.
\end{enumerate}

This design ensures that factual claims are derived from verified institutional documents rather than generated from the model's parametric memory, directly reducing hallucination risk as discussed in \secref{sec:llm_foundations}.

\subsection{Room Navigation Lookup}
\label{sec:room_lookup}

A common receptionist query type is asking for directions to a specific room.
Rather than relying on vector search for this structured data, the system uses a deterministic lookup mechanism embedded in the \texttt{query\_search} tool.

Before forwarding a query to Weaviate, the tool checks whether the input matches a room-related pattern (a room number combined with directional keywords such as ``room'', ``where'', or ``direction'').
If a match is found, the query is routed to a lookup table of building rooms organized by floor, each entry containing natural-language walking directions from the robot's position.
The result is returned directly without a vector search, providing immediate responses for navigation queries.

This routing was originally implemented as a separate tool.
However, testing revealed that the local LLM's tool-calling reliability degraded with additional tool schemas (see \secref{sec:local_mode}), so the functionality was merged into \texttt{query\_search} to maintain a minimal two-tool interface.


% ============================================================
\section{Non-Verbal Behavior}
\label{sec:nonverbal}

To create a more natural interaction experience, the system accompanies spoken dialogue with physical gestures on the Pepper robot.

\subsection{Gesture and Animation System}
\label{sec:animations}

Pepper's gestures are triggered through the \texttt{play\_animation} tool, which the LLM calls as part of generating each response.
The tool interface accepts a single argument: a semantic label describing the intended gesture category.

The system defines nine animation groups, each mapping a semantic label to a set of concrete Pepper animation sequences:

\begin{table}[h]
\centering
\begin{tabular}{l l c}
\hline
\textbf{Group} & \textbf{Description} & \textbf{Variants} \\
\hline
\texttt{greeting}   & Waving, hello gestures     & 9 \\
\texttt{bow}        & Short bow, acknowledgement  & 3 \\
\texttt{explain}    & Explanatory hand gestures   & 8 \\
\texttt{happy}      & Cheerful, positive gestures & 4 \\
\texttt{thinking}   & Contemplative poses         & 8 \\
\texttt{dont\_know} & Shrug, uncertainty          & 6 \\
\texttt{excited}    & Enthusiastic gestures       & 3 \\
\texttt{interested} & Leaning in, attentive       & 2 \\
\texttt{surprised}  & Surprise reactions          & 4 \\
\hline
\end{tabular}
\caption{Pepper animation groups. Each group maps a semantic label to multiple animation variants. The system selects a random variant on each invocation to avoid repetitive motion.}
\label{tab:animation_groups}
\end{table}

When the tool is invoked, the input is normalized and resolved through a chain: direct group name match, then natural-language alias lookup (e.g., ``hello'' resolves to \texttt{greeting}, ``goodbye'' to \texttt{bow}), then direct animation key match.
A random variant is selected from the matched group to ensure varied, non-repetitive motion across interactions.

The resolved animation is dispatched to the robot bridge, which mutes the robot's audio to prevent sound interference, executes the animation via NAOqi, and restores audio --- all asynchronously.
The HTTP response is returned immediately, ensuring that animation dispatch does not block the dialogue turn or starve the agent's WebRTC connection.

The system prompt instructs the LLM to call the animation tool once per reply, selecting a gesture appropriate to the emotional tone of the response.
In cloud mode, the available groups are restricted to six core categories for reliability; in local mode, all nine groups are accessible.

\subsection{Tablet Interface}
\label{sec:tablet}

Pepper's chest-mounted tablet serves as a secondary visual output channel.
The robot bridge exposes endpoints for displaying inline HTML content (with configurable styling), rendering a conversation history as a chat interface, and clearing the display.
Content is delivered to the tablet through NAOqi's tablet service, which renders HTML locally on the tablet's built-in browser.

The tablet is currently used to display status information and conversation context.
Full integration with the dialogue system (e.g., showing retrieved search results or navigation maps) is planned as future work.


% ============================================================
\section{Session Management and Orchestration}
\label{sec:orchestration}

\subsection{Session Lifecycle}
\label{sec:session_lifecycle}

The session manager is the central orchestration service.
It manages the lifecycle of interaction sessions, which proceed through the following states:

\begin{enumerate}
    \item \textbf{Idle:} No active conversation.
    A voice agent is connected to the LiveKit room as a warm (preloaded) participant but is not engaged in dialogue.
    The session manager monitors for incoming user connections.

    \item \textbf{Active:} A user is present in the room and dialogue is in progress.
    The session manager sends an activation signal to the voice agent via a LiveKit data channel, prompting it to begin interaction.

    \item \textbf{Reset:} After the user disconnects or an idle timeout elapses (60 seconds of silence by default), the session manager sends a reset signal.
    The voice agent clears its conversation history in place --- without tearing down the WebRTC connection or unloading models --- and returns to the idle state, ready for the next visitor.
\end{enumerate}

The warm agent model is a deliberate design choice to minimize latency between visitor arrival and the first response.
Loading speech processing models (VAD, STT, TTS) takes several seconds on the target hardware; keeping the agent process alive across sessions amortizes this startup cost.
Temporary user disconnections (e.g., browser refresh, brief network interruption) do not destroy the agent's state --- only an explicit shutdown signal from the session manager terminates the process.

The session manager also runs a periodic monitoring loop that detects missing agent participants.
If the voice agent disappears from the room due to a crash or network failure, a grace period is applied before automatic redispatch, preventing false positives during brief reconnection windows.

\subsection{Agent Dispatch and Mode Switching}
\label{sec:agent_dispatch}

Two voice agents can be registered with LiveKit simultaneously, each under a unique name: one for cloud mode (running on the RPi) and one for local mode (running on the GPU server).
The session manager dispatches to the correct agent based on the currently selected mode, passing the mode designation as metadata in the dispatch request.

When the mode is changed via the Operator Panel or REST API, the session manager sends a shutdown signal to the currently active agent.
The agent exits its persistent event loop, freeing the worker process, and the session manager dispatches a new warm agent of the requested mode into the same room.
This process takes several seconds (primarily due to model loading in local mode) but requires no manual intervention beyond the mode selection.

\subsection{Audio Delivery to the Robot}
\label{sec:audio_delivery}

The voice agent generates audio within the LiveKit room, but Pepper's speakers are not directly accessible from LiveKit.
The audio bridge service closes this gap by acting as a protocol translator between WebRTC and Pepper's native audio interface.

The audio bridge joins the LiveKit room as a non-agent participant and subscribes to the voice agent's audio track.
Incoming audio frames are resampled to 16\,kHz mono (matching Pepper's expected input format) and attenuated to prevent distortion on the robot's speakers.
The processed frames are sent over a TCP connection to the robot bridge, which buffers and plays them through NAOqi's audio device service.

The audio bridge is designed to be resilient: it starts independently of the robot bridge, silently dropping audio frames if the TCP connection is unavailable, and automatically reconnects when the bridge becomes reachable.
When a new agent joins the room after a mode switch, the audio bridge detects the change and seamlessly redirects its subscription to the new agent's audio track.


% ============================================================
\section{Implementation}
\label{sec:implementation}

\subsection{Software Stack}
\label{sec:software_stack}

\tabref{tab:software_stack} summarizes the key technologies used in the system.
All components are implemented in Python 3.12.
The \texttt{qi} library for Pepper communication was self-built from source for ARM64 on the Raspberry Pi, as no pre-built package is available for this platform.

\begin{table}[h]
\centering
\begin{tabular}{l l}
\hline
\textbf{Component} & \textbf{Technology} \\
\hline
WebRTC / Rooms      & LiveKit Server + LiveKit Agents SDK \\
LLM (cloud)         & OpenAI Realtime API (\texttt{gpt-4o-mini-realtime}) \\
LLM (local)         & vLLM + Qwen 2.5 7B Instruct (AWQ) \\
STT (local)         & Faster Whisper (CTranslate2) \\
TTS (local)         & Piper (VITS-based, ONNX) \\
VAD                 & Silero VAD \\
Vector Database     & Weaviate \\
Embeddings          & OpenAI \texttt{text-embedding-3-large} \\
Robot SDK           & libqi / NAOqi \\
Containerization    & Docker Compose \\
\hline
\end{tabular}
\caption{Software stack overview.}
\label{tab:software_stack}
\end{table}

\subsection{Deployment and Containerization}
\label{sec:deployment}

All services on the RPi are defined in a single Docker Compose configuration.
A shared base image provides Python 3.12 with the \texttt{uv} package manager; each service mounts the project source code as a bind volume, eliminating the need to rebuild container images after code changes.

The robot bridge container requires host network access to reach Pepper's NAOqi service, along with volume mounts for the self-built \texttt{qi} bindings and their native library dependencies.
Other containers communicate through Docker's internal networking.

Optional services --- such as a microphone client for RPi-based audio input, SSH tunnel containers for GPU server connectivity, and a development console --- are organized into Docker Compose profiles, allowing flexible startup configurations depending on the deployment scenario.

\subsection{Reproducibility}
\label{sec:reproducibility}

To reproduce the system, the following are required: a Raspberry Pi 5, a SoftBank Pepper robot on the same network, API keys for OpenAI and LiveKit, and (for local mode) access to a GPU server with vLLM.
The system is started with a single Docker Compose command; the Operator Panel then provides a web interface for session monitoring, mode switching, and transcript viewing.
Full setup instructions and configuration details are provided in the project repository.
